\section{Time penalties and H\"anni's collapse}\label{sec:time}
% =====================================================================
% Section 6: the collapse of axiom induction into function induction, what templates do to it, and which
% time or length penalty prices it. Sources (final records only):
%   research/tracks/model/notes-final.md   §6 (Prop 6.0 - Conj 6.15, the closing statement), §1.3, §1.5,
%       Rem 1.10, §9.1 (H6), §9.3, §10 items 12-15, §12 (verification log: M2, m1); checks/c12_l1_size.out
%   research/tracks/model/referee.md        (M2, m1, Q2); referee_code/r2_l1_size.out
%   research/tracks/universal/notes-final.md Prop U2t (§2.2), the remark after Prop U14 (§8)
%   research/tracks/pa/notes-final.md       Prop 4.3, §3.2, §3.8, F9; pa referee.md missed question 3
%   research/tracks/experiments/notes-final.md §1.3, §9 (E7); code/results/e7_prior.md
%   research/prior/hanni-*.md (quotes)
% Proofs and computations: sections/app-time.tex.
% Local deviations from NOTATION.md §9: the binary numeral of a word w is \num{w} (not nu(w)), because nu(pi)
% counts grammar nodes in the same section; the exponential rate of model Prop 6.11 is zeta (the notes'
% theta), because theta is a substitution (as in app-model.tex).
% This section is the canonical home of Haenni's collapse argument and its quotations (DECISIONS §4).
% =====================================================================

H\"anni's working notes \citep{hanni2026notes}\footnote{File \texttt{research/prior/hanni-solomonoff-axiom-induction.md}; working notes, quoted verbatim.} compare axiom induction, which weighs axiom systems that prove the labelled data, with Solomonoff function induction, which weighs programs that label statements directly. He argues that the two are equivalent for consistent assigners: an assigner $T$ becomes the schema ``\texttt{T(quoted-phi) = accept => phi}'', where ``being an axiom of the right form is in fact decidable'', and the result is consistent because ``a model of the phis alone would be a model of the added sentences also''. So unrestricted axiom induction collapses into function induction, and he considers restricting axioms to ``second-order substitution schemas of a specific kind''. The question adds a time penalty: ``it should not take too long to generate axioms / to check whether a statement is an axiom''. The result (proofs in \cref{app:time}):

\begin{remark}[What the time penalty does]\status{(1), (2), (4) proved; (3) lower bound proved, upper bound proof sketch}\src{model \S6.5}\label{rem:time:summary}
(1) Penalties on the time to check or generate axioms (Kt, speed prior) do not block the collapse: Craig's sets $A^C_f$ have membership decidable in polynomial time (\cref{prop:time:cheap}), like finite $\DT$ theories, so such a penalty charges them as it charges template theories; no time-penalised prior is defined here, so no bound on how much it changes \cref{thm:time:equiv} is claimed. (2) Templates block H\"anni's schema (its acceptance part, \cref{prop:time:notemplate}) but, over $\PA$, not the collapse for $\Sigma_n$-sound assigners on $\Sigma_n$ sentences: one ground reflection sentence reproduces them at a cost linear in $|f|$ (\cref{prop:time:collapse}); for merely consistent assigners that sentence can be inconsistent with $\PA$, and whether templates block the collapse for them is open. (3) Derivation size in written symbols prices the assigner: one decidable $X$ for all finite $\DT$ theories, and one for each $e$ for all axiom sets with membership in time $O(n^e)$, forces every such theory consistent with $\Gamma_{f_X}$ to use derivations of symbol size above $s(m)$ for infinitely many lengths $m$, i.e.\ at least a fixed polynomial root of the assigner's nondeterministic time (\cref{thm:time:ntime,cor:time:hard}); $\Lsig$ then charges at least $\kappa s(m)/4-\ln(Z_T/Z^\sigma_T)$ nats and the graded score at least $\kappa s(m)$ bits ($\kappa s(m)+\log_2Z_T$ after normalisation) on those data (\cref{prop:time:sigma}). The collapse constructions through $A^C_f$ and the two-sorted $A_f$ pay at most a polynomial of the deterministic running time (proof sketch, \cref{rem:time:upper}); for $\rho_{f,n}$ no upper bound is established. (4) Plain $\Lone$ is proved to charge at least $\zeta s(m)-\ln(C/Z_T)$ nats when the assigner's language lies outside $\NTIME(2^{(d+1)s})$, about the logarithm of its nondeterministic time (\cref{cor:time:log}); a polynomial charge is \cref{conj:time:poly}. Function induction's total log loss on $f_X$'s labels is at most $|f_X|$ bits (when $B\cup\Gamma_{f_X}$ is consistent).
\end{remark}

% ---------------------------------------------------------------------
\subsection{The two inductors}\label{sec:time:inducers}

Fix a countable first-order language $L$ with sentences $\Sent_L$, a universal prefix machine $U$ (program $p$ has weight $2^{-|p|}$), and a decidable background $B\subseteq\Sent_L$, possibly empty, that is part of every axiom hypothesis. Data are labelled sentences $D=((\varphi_1,b_1),\dots,(\varphi_n,b_n))$, $b_i\in\{1,0\}$ (true, false).

\begin{definition}[Axiom and function induction]\src{model \S6.1}\label{def:time:inducers}
\begin{itemize}[leftmargin=*,itemsep=1pt]
\item $\AI$: a hypothesis is a program $p$ that halts on every input and decides a set $A_p\subseteq\Sent_L$. It is \emph{compatible} with $D$ if $B\cup A_p$ is consistent and proves $\varphi_i$ when $b_i=1$ and $\neg\varphi_i$ when $b_i=0$. $W_{\AI}(D):=\sum_{p\text{ compatible}}2^{-|p|}$.
\item $\FIcons$: a hypothesis is a program $f$ computing a partial map $\Sent_L\to\{\mathrm{acc},\mathrm{rej}\}$ with $B\cup\Gamma_f$ consistent, where $\Gamma_f:=\{\varphi:f(\varphi)=\mathrm{acc}\}\cup\{\neg\varphi:f(\varphi)=\mathrm{rej}\}$; it is compatible if it gives each $\varphi_i$ the label $b_i$; $W_{\mathrm{FI}}(D):=\sum_{f\text{ compatible}}2^{-|f|}$. $\FIall$ is the same over all assigners.
\item Semimeasure convention: $P_X(b_1\dots b_n\mid\varphi_1\dots\varphi_n):=W_X(D_n)/W_X(\emptyset)$, and log loss is $-\log_2$ of it. Hypotheses that leave $\varphi$ undecided drop out, so $W(D+(\varphi,1))+W(D+(\varphi,0))\le W(D)$ (H\"anni: ``maybe we should speak of a semimeasure here instead of renormalizing''). $W_X(D)$ depends only on the set $D$.
\end{itemize}
\end{definition}

\noindent Function induction over all assigners is not the fair comparison: on a consistent labelling of fair coin flips $\FIall$ has bounded log loss while $\AI$ and $\FIcons$ lose at least $n$ bits in expectation (\cref{prop:time:fiall}), as H\"anni argued informally.

\begin{theorem}[H\"anni's equivalence, repaired]\status{proved; Craig's theorem known}\src{model Thm 6.1}\label{thm:time:equiv}
There is a constant $c$, depending only on $U$, the coding of syntax and the decider for $B$, such that $2^{-c}W_{\AI}(D)\le W_{\mathrm{FI}}(D)\le2^{c}W_{\AI}(D)$ for every finite labelled $D$. Hence for every reveal order and labelling the semimeasure log losses of $\AI$ and $\FIcons$ differ by at most $2c$ bits.
\end{theorem}

\begin{proof}[Proof idea]
$\AI\to\mathrm{FI}$ is H\"anni's argument: map $p$ to ``search for proofs of $\varphi$ and of $\neg\varphi$ from $B\cup A_p$''. For $\mathrm{FI}\to\AI$ his schema is replaced by Craig's trick \citep{craig1953axiomatizability}: $f$ becomes a decider of
$A^C_f:=\{\varphi^{\wedge(k+1)}:f(\varphi)=\mathrm{acc}\text{ in exactly }k\text{ steps}\}\cup\{(\neg\varphi)^{\wedge(k+1)}:f(\varphi)=\mathrm{rej}\text{ in exactly }k\text{ steps}\}$, $\psi^{\wedge m}$ the right-nested $m$-fold conjunction. $A^C_f$ is decidable and has the consequences of $\Gamma_f$; its decider is $f$ plus a fixed wrapper (\cref{app:time:inducers}).
\end{proof}

The constant equivalence is proved only in the semimeasure convention; renormalising at each step, the argument gives only $2c$ bits per step (\cref{rem:time:convention}).

% ---------------------------------------------------------------------
\subsection{H\"anni's schema}\label{sec:time:schema}

Let $\Accf{f}(x)$ be a $\Sigma_1$ formula of arithmetic for ``$f$ halts on input $x$ with output $\mathrm{acc}$'' (via Kleene's $T$ predicate), and $\Rejf{f}(x)$ likewise. Both are written without $<$ (bounded quantifiers as $\exists z\,(z+x=y)$), so that $\Sigma_1$-completeness of $\Q$, which in $L_A$ holds for $<$-free $\Sigma_1$ sentences, applies. H\"anni's schema is
$A_f:=\{\Accf{f}(\gn{\varphi})\to\varphi:\varphi\in\Sent_L\}\cup\{\Rejf{f}(\gn{\varphi})\to\neg\varphi:\varphi\in\Sent_L\}$.

\begin{proposition}[Two sorts: his argument works]\status{proved; $\Sigma_1$-completeness of $\Q$ for $<$-free sentences known}\src{model Prop 6.2}\label{prop:time:twosorted}
Let $L$-sentences live in one sort and arithmetic in another, with $B:=\Q$ on the second. If $\Gamma_f$ has a model, then $B\cup A_f$ is consistent and $\Cn(B\cup A_f)\cap\Sent_L=\Cn(\Gamma_f)$.
\end{proposition}

\begin{proposition}[One sort over $\PA$]\status{proved; H\"anni's argument refuted in this reading}\src{model Prop 6.3}\label{prop:time:single}
Assume $\Con(\PA)$ and the standard formalisation of computations. Let $L=L_A$, $B=\PA$, and let $f$ accept $\gn{\neg\Con(\PA)}$, search on input $\gn{\Con(\PA)}$ for a $\PA$-proof of $\bot$ and accept if it finds one, and loop otherwise. Then $\Gamma_f=\{\neg\Con(\PA)\}$ is consistent with $\PA$; $\PA\cup A_f$ is not.
\end{proposition}

His consistency argument assumes that the antecedents are evaluated correctly; in a model of $\PA+\neg\Con(\PA)$, $\Accf{f}(\gn{\Con(\PA)})$ holds via a nonstandard proof of $\bot$. So his \emph{conclusion}, the equivalence, stands (\cref{thm:time:equiv}); his \emph{argument} needs two sorts or a sound $f$.

% ---------------------------------------------------------------------
\subsection{What templates do}\label{sec:time:templates}

\begin{proposition}[No non-ground template in the collapse schema]\status{proved}\src{model Prop 6.4; pa Prop 4.3}\label{prop:time:notemplate}
Write Gödel numbers as unary numerals $S^k0$ (or any numerals in which some term symbol, here $+$, never occurs), and let $C_f:=\{\Accf{f}(\gn{\varphi})\to\varphi:\varphi\in\Sent_L\}$, the acceptance part of $A_f$, single- or two-sorted. Every $\DT$ template $\tau$ with $\inst(\tau)\subseteq C_f$ is ground. So $C_f$ is not $\bigcup\inst(T)$ for any finite set $T$ of $\DT$ templates, and such a set with instances in $C_f$ covers finitely many of its members. Likewise for the reflection schema $\Refl_T=\{\Prv_T(\gn{\varphi})\to\varphi\}$ (\cref{prop:pa:refl}).
\end{proposition}

\noindent The rejection part, and templates whose instances mix the two parts, are not treated here (\cref{app:ver:open}). For the acceptance part, H\"anni's restriction blocks his construction: inside $\DT$ it costs one ground template per sentence, which is memorisation. But one ground sentence can carry the collapse over $\PA$.

\begin{proposition}[One reflection sentence carries the collapse]\status{proved, given partial truth predicates (known)}\src{model Prop 6.5}\label{prop:time:collapse}
For $n\ge1$ let $\Tr_n(x)$ be a $\Sigma_n$ partial truth predicate, $\PA\vdash\varphi\leftrightarrow\Tr_n(\gn{\varphi})$ for every $\Sigma_n$ sentence $\varphi$ (known; \citealp[Ch.~I]{hajek1993metamathematics}; \citealp[Ch.~9]{kaye1991models}; chapter numbers not checked), and $\Sent_{\Sigma_n}(x)$ a $\Delta_1$ formula for ``$x$ codes a $\Sigma_n$ sentence''. The single sentence (ground $\DT$ template)
\[
\rho_{f,n}:=\forall x\big[(\Sent_{\Sigma_n}(x)\wedge\Accf{f}(x))\to\Tr_n(x)\big]\wedge\forall x\big[(\Sent_{\Sigma_n}(x)\wedge\Rejf{f}(x))\to\neg\Tr_n(x)\big]
\]
has $\ell(\rho_{f,n})\le a|f|+b_n$ for constants $a,b_n$ (binary numeral for $f$'s index), and:
(a) $\PA+\rho_{f,n}$ proves every $\Sigma_n$ sentence that $f$ accepts and the negation of every one it rejects;
(b) if $f$ is $\Sigma_n$-sound (it accepts only true and rejects only false $\Sigma_n$ sentences), $\PA+\rho_{f,n}$ is true in $\N$.
\end{proposition}

So $\PA$'s templates plus one ground sentence reproduce any $\Sigma_n$-sound assigner on $\Sigma_n$ sentences, at a cost linear in $|f|$ (not $|f|+c$); for merely consistent assigners $\rho_{f,n}$ can be inconsistent with $\PA$ (as in \cref{prop:time:single}).

% ---------------------------------------------------------------------
\subsection{Time penalties}\label{sec:time:penalties}

\begin{proposition}[Membership is cheap in both collapse sets]\status{proved}\src{model Prop 6.6}\label{prop:time:cheap}
(a) With binary numerals, membership of $\chi$ in $A_f$ is decidable in time polynomial in $|\chi|$ without running $f$: parse $\chi$ as $\Accf{f}(t)\to\varphi$ or $\Rejf{f}(t)\to\neg\varphi$, compute $\gn{\varphi}$, compare with $t$. With unary numerals likewise, since then $|\chi|\ge\gn{\varphi}$.
(b) Membership in $A^C_f$ is decidable in time polynomial in $|\chi|$: at most $|\chi|$ decompositions $\chi=\psi^{\wedge m}$, each running $f$ for $m-1<|\chi|$ steps.
\end{proposition}

So a penalty in the style of Levin's $\mathrm{Kt}$ or the speed prior on the time to \emph{check} axiomhood puts these sets in the same polynomial class as $\DT$ templates (\cref{rem:model:priors}(iii)): it does not price the collapse. Penalising the time to \emph{generate} axioms fails too: $\varphi^{\wedge(k+1)}\in A^C_f$ comes from $k$ steps of $f$ and is longer than $k$.

The computation of $f$ moves into the proofs ($A_f$: a derivation must prove $\Accf{f}(\gn{\varphi})$) or into the axioms' length ($A^C_f$), both counted by derivation size in symbols; no theory with polynomial-time membership escapes this. Let $L\supseteq\{0,S,+,\cdot,R\}$, $R$ unary. For $X\subseteq\{0,1\}^*$ put $\varphi_w:=R(\num{w})$, $\num{w}$ a binary numeral term, and let $f_X$ accept $\varphi_w$ iff $w\in X$ and reject it otherwise; $\Gamma_{f_X}$ is consistent (interpret $R$ as $X$).

\begin{theorem}[Derivation size bounds the assigner's nondeterministic time]\status{proved}\src{model Thm 6.7}\label{thm:time:ntime}
Let $T$ be any axiom set with membership decidable in time $O(n^e)$ and $T\cup\Gamma_{f_X}$ consistent. If every $\varphi_w$, $w\in X$, has a $\CK$-derivation from $T$ of symbol size at most $\ell(|w|)$, with $\ell(m)\ge m$ time-constructible, then $X\in\NTIME(\ell(m)^{c_e})$, $c_e$ depending only on $e$ and $\CK$'s checking overhead.
\end{theorem}

\begin{proof}[Proof idea]
Guess a derivation of size at most $\ell(|w|)$ and check it. If $w\notin X$, then $\neg\varphi_w\in\Gamma_{f_X}$, so consistency rules out a derivation (\cref{app:time:penalties}).
\end{proof}

\begin{corollary}[Hard assigners force long derivations, for every theory]\status{proved, given the nondeterministic time hierarchy theorem (known)}\src{model Cor 6.8 (referee m1)}\label{cor:time:hard}
For finite $\DT$ theories membership takes time $O(|T|\cdot|\chi|)$ (AS Thm~2.5), so $e=1$ and $c_1$ is one constant for all of them.
(a) Let $s(m)\ge m$ be time-constructible with $s(m+1)^{c_1}=o(s(m)^{c_1+1})$ (every polynomial). There is a decidable $X$ such that every finite $\DT$ theory $T$ with $T\cup\Gamma_{f_X}$ consistent has, for infinitely many $m$, some $w\in X$ of length $m$ with no $\CK$-derivation of $\varphi_w$ from $T$ of symbol size $\le s(m)$.
(b) For each $e$ the same holds, with $X$ depending on $e$, for every axiom set with membership in time $O(n^e)$; a single $X$ for all $e$ exists under a growth condition \status{proof sketch} (\cref{app:time:penalties}).
\end{corollary}

\noindent The hierarchy theorem is cited from \citet{cook1972hierarchy,seiferas1978separating,zak1983turing}, unchecked.

% ---------------------------------------------------------------------
\subsection{What the derivation likelihoods charge}\label{sec:time:lik}

\Cref{thm:time:ntime} bounds the \emph{symbol} size of derivations. The code length $-\ln\Pr(\pi)$ of an $\Lone$ tree counts its grammar choices instead (\cref{def:model:lone}), and the two can differ exponentially.

\begin{remark}[Refuted: plain $\Lone$ charges symbol size]\status{refuted; counterexample proved and computed}\src{model Prop 6.9}\label{prop:time:lonesize}
Under $\Lone$ with subcritical parameters, $-\ln P_T(\varphi)\ge\kappa\,\ell_{\min}(\varphi)-\ln(C/Z_T)$, $\ell_{\min}$ the least symbol size of a derivation, fails (an earlier sketch claimed it; \cref{app:ver:corrections}). From $\forall x\,(x=x)$, repeat $k$ times ``$\forall$E with $t=p+p$, then Gen on $p$'', and close with one $\forall$E: the conclusion has size $2^{k+3}-1$, while $-\ln\Pr(\text{tree})=12.77+12.21k$ nats ($12.206$ per round) and $-\ln P_T\le-\ln\Pr(\text{tree})$ (\cref{app:time:lik}, \cref{tab:time:c12}).
\end{remark}

$\forall$E and Gen substitute into every occurrence of a variable at the price of one choice, so \cref{thm:time:ntime,cor:time:hard} transfer to $\Lsig$ and to H\"anni's graded score, not to plain $\Lone$. Let $\nu(\pi)$ count the grammar nodes of a tree $\pi$ (derivation, $\Qg$ and parameter-chain nodes) and $\nu_{\min}(s)$ be the least $\nu$ of a valid tree concluding $s$. $\Lone$'s code length counts grammar choices: a tree with $\nu$ grammar nodes from a finite $T$ has conclusions of symbol size at most exponential in $\nu$, and $\Lone$ with a uniformly subcritical grammar and a geometric parameter law charges at least $\zeta\nu_{\min}(\varphi)-\ln(C/Z_T)$ nats for $\varphi$; so if each $\varphi_w$, $w\in X$, has a tree with at most $\ell(|w|)$ nodes from a finite $\DT$ theory consistent with $\Gamma_{f_X}$, then $X\in\NTIME(2^{(d+1)\ell(m)})$, $d$ independent of $T$ (\cref{lem:time:symexp,prop:time:codelength,thm:time:log} in \cref{app:time:lik}).

\begin{corollary}[Hard assigners make plain $\Lone$ pay, logarithmically]\status{proved, given the hierarchy theorem (known)}\src{model Cor 6.13}\label{cor:time:log}
Let $s(m)\ge m$ be time-constructible with $s(m)^2-(d+1)s(m+1)\to\infty$ (every polynomial; $2^m$). There is a decidable $X$ such that for every finite $\DT$ theory $T$ with $T\cup\Gamma_{f_X}$ consistent, for infinitely many $m$ some $w\in X$ of length $m$ has $\nu_{\min}(\varphi_w)>s(m)$; under the hypotheses of \cref{prop:time:codelength}, $-\ln P_T(\varphi_w)\ge\zeta\,s(m)-\ln(C/Z_T)$ on those data.
\end{corollary}

Symbol-size penalties, by contrast, charge a polynomial root of the nondeterministic time:

\begin{proposition}[Symbol-size penalties charge polynomially]\status{(a) proved; (b), (c) proved, given the hierarchy theorem}\src{model Prop 6.14}\label{prop:time:sigma}
(a) Under $\Lsig$ (\cref{def:model:variants}), $P^\sigma_T(\varphi)\le e^{-\kappa\,\ell^\sigma_{\min}(\varphi)}Z_T/Z^\sigma_T$, with $\ell^\sigma_{\min}(\varphi)$ the least total symbol size of a valid $\Lone$ tree concluding $\varphi$.
(b) A valid $\Lone$ tree of total symbol size $\ell$ yields a $\CK$-derivation of symbol size at most $4\ell$. So with $X$ from \cref{cor:time:hard}(a), for every finite $\DT$ theory $T$ consistent with $\Gamma_{f_X}$, infinitely often $-\ln P^\sigma_T(\varphi_w)\ge\kappa s(m)/4-\ln(Z_T/Z^\sigma_T)$.
(c) H\"anni's graded score with $g(\ell)=2^{-\kappa\ell}$ on explicit $\CK$-derivations, in a code with at least one bit per written symbol (\cref{rem:model:graded}), charges at least $\kappa s(m)$ bits on the same data ($\kappa s(m)+\log_2Z_T$ normalised).
\end{proposition}

\begin{conjecture}[Plain $\Lone$ charges polynomially too]\status{conjecture}\src{model Conj 6.15}\label{conj:time:poly}
\Cref{thm:time:log} holds with $\NTIME(\ell(m)^c)$, for a constant $c$, in place of $\NTIME(2^{(d+1)\ell(m)})$.
\end{conjecture}

The missing step is checking trees on shared representations, with de Bruijn indices and Gen under sharing (\cref{app:time:lik}). Unless \cref{conj:time:poly} is proved, the question's derivation-length prior should therefore measure length in written symbols ($\Lsig$, or the graded score normalised as in \cref{rem:model:graded}). Whether axiom induction with a symbol-size penalty is equivalent, up to polynomial factors, to function induction over consistent assigners with short certificates is open.

% ---------------------------------------------------------------------
\subsection{Links to the other sections}\label{sec:time:summary}

\noindent Links to \cref{sec:univ,sec:pa} are in \cref{rem:time:links}: no penalty monotone in the number of derivation steps reverses \cref{thm:univ:B}; a Kt-style penalty charges only $O(\ln k)$ for $0+S^k0=S^k0$ on the equational fragment; and a time-bounded likelihood favours making theorems with long proofs axioms.

\begin{remark}[Prior variants on the E2 data]\status{computed}\src{experiments E7}\label{rem:time:e7}
On the E2 data (25 seeds), a membership-time factor $(1+|T|)^{-\tau}$ shifts the log-odds of $T^*$ against a hand competitor by at most $\log_2(78/2)\approx5.3$ bits at $\tau=1$; it changes nothing visible from $n=64$ on and at $n\le32$ favours small unsound lumps, slightly at $\tau=1$ (unsound MAP at $n=8$ in 9 instead of 7 seeds) and markedly at $\tau=4$ (mean unsound mass at $n=8$ $0.372$ against $0.256$; a false probe accepted in 5 instead of 2 seeds). A steeper simplicity penalty $\lambda=2$ makes the MAP at $n=8$ unsound in 22 seeds ($\lambda=1$: 7) but removes the false spare slot faster ($2\cdot10^{-4}$ against $0.004$ unsound mass at $n=512$); $\lambda=\tfrac12$ (proper only on finite pools) the reverse: neither removes both failures (\cref{tab:exp:e7}).
\end{remark}

\begin{conjecture}[A time-bounded analogue]\status{conjecture; not written out}\src{model \S9.3}\label{conj:time:polytime}
H\"anni's note on time-bounded Solomonoff induction\footnote{File \texttt{research/prior/hanni-polytime-solomonoff.md}; working notes, in which he says he has not ``checked all the technical claims completely carefully''.} gives, for each polynomial $p$, a predictor whose log loss exceeds that of any ``prediction algorithm $P$ with runtime $O(p(n))$'' by ``at most a constant $c_P$''. By his argument, the derivation-bounded $\Ltwo$ inducer with a growing depth budget plausibly has constant regret against theories whose data have derivations within the budget.
\end{conjecture}
